% !PRINTED 31
\weberpage{31}
%<<CONT
gesetzt wird, das Differential (5) in
\weq{8}{\frac{1}{M}\,\frac{\xi\,d\eta - \eta\,d\xi}{\sqrt{\varphi(\xi,\eta)}}}
übergeht. Es läßt sich nun nachweisen, daß, sobald die Bedingung (6)
erfüllt ist, $H$ durch $T$ teilbar, und also, da der Grad beider
Funktionen derselbe ist, $M$ eine Konstante wird. Diese Konstante heißt
der \gesperrt{Multiplikator} der Transformation.

Bemerken wir nämlich, daß, da $f(x,y)$ keinen quadratischen Faktor
enthält, die beiden Funktionen
\[
\frac{\partial f}{\partial x}, \quad \frac{\partial f}{\partial y}
\]
keinen gemeinschaftlichen Teiler haben, und daß infolgedessen, weil $U$
und $V$ teilerfremd sind, auch
\[
\frac{\partial f}{\partial U}, \quad \frac{\partial f}{\partial V}
\]
keinen gemeinschaftlichen Teiler (in Beziehung auf $\xi,\eta$) haben, so
lassen sich zwei Funktionen $\alpha,\beta$ von $\xi,\eta$ so bestimmen,
daß
\[
\alpha\,\frac{\partial f}{\partial U} + \beta\,\frac{\partial f}{\partial V}
\]
zu einer beliebig gegebenen Funktion $T$ teilerfremd ist. Man kann z.\,B.
$\alpha$ teilerfremd zu $T$ und $\beta$ durch die in $\partial f/\partial U$
nicht aufgehenden Teiler von $T$ teilbar, dagegen durch die
gemeinschaftlichen Teiler von $T$ und $\partial f/\partial U$ unteilbar
annehmen. Wenn wir nun die Gleichung (6), die wir als erfüllt
voraussetzen, nach $\xi,\eta$ differentiieren, so folgt:
\[
\begin{aligned}
\frac{\partial f}{\partial U}\frac{\partial U}{\partial \xi} + \frac{\partial f}{\partial V}\frac{\partial V}{\partial \xi} &= TX,\\
\frac{\partial f}{\partial U}\frac{\partial U}{\partial \eta} + \frac{\partial f}{\partial V}\frac{\partial V}{\partial \eta} &= TY,
\end{aligned}
\]
und daraus durch Auflösung in bezug auf $\partial f/\partial U$,
$\partial f/\partial V$:
\[
H\left(\alpha\,\frac{\partial f}{\partial U} + \beta\,\frac{\partial f}{\partial V}\right) = TZ,
\]
worin $X,Y,Z$ ganze homogene Funktionen von $\xi,\eta$ sind. Aus der
letzten Gleichung aber schließt man, daß $H$ durch $T$ teilbar sein muß.
